\section{Reproducibility And Artifact Availability}

The paper exposes four evidence surfaces rather than a single undifferentiated bundle: the main evidence surface, an appendix-only supplementary diagnostic, a manual audit bundle, and benchmark-specific reliability notes. Each surface is frozen locally and is named explicitly so that readers can verify the headline claims without rerunning the experiments.

\paragraph{Main evidence surface.}
The reportable external evidence surface is the registry-listed set of official bundles under \texttt{artifacts/tmlr\_official/}. These bundles cover PopQA recurring memory, FreshQA public main and sequence, MQuAKE, UniEdit, and the cross-benchmark statistics bundle. For any official run root, the first inspection files are \texttt{manifest.json}, \texttt{aggregate\_table.md}, \texttt{run\_rows.jsonl}, \texttt{significance.json}, and \texttt{bootstrap.json}; MQuAKE and UniEdit additionally expose benchmark-specific aggregates through \texttt{benchmark\_metric\_aggregates.json}. The cross-benchmark law is summarized in \texttt{artifacts/tmlr\_official/official\_statistics\_v1/}.

\paragraph{Appendix-only supplementary diagnostic.}
The only supplementary diagnostic kept outside the main external evidence surface is \texttt{artifacts/freshness\_v1/}. It is a deterministic local controlled-update benchmark used as extra arXiv context. It remains useful for understanding stable update behavior, but it is not part of the official external benchmark surface and should not be read as overriding the PopQA, FreshQA public, MQuAKE, or UniEdit results.

\paragraph{Manual audit.}
The public benchmark trust story is grounded in \texttt{artifacts/tmlr\_official/freshqa\_leakage\_audit\_manual\_v1/}. The main inspection files are \texttt{summary.json}, \texttt{manual\_audit\_100.md}, \texttt{manual\_audit\_100.json}, and \texttt{selection\_manifest.json}. These files expose the bounded trust conclusion directly: no direct leakage-risk cases were found, but many audited rows remain suspect.

\paragraph{Reliability notes.}
Benchmark-construction and field-mapping caveats are separated from the main tables. The MQuAKE reliability note lives in \texttt{artifacts/tmlr\_official/mquake\_reliability\_note\_v1/}, and the UniEdit snapshot note lives in \texttt{artifacts/tmlr\_official/uniedit\_reliability\_note\_v1/}. These notes document how the benchmark inputs were frozen locally and how benchmark-specific metrics should be interpreted.

\paragraph{How to verify.}
\begin{itemize}[leftmargin=1.5em]
\item Open \texttt{paper\_run\_manifest\_index.json} to see which artifact roots are reportable for the paper.
\item Trace each headline claim through \texttt{paper/shared/claims\_to\_artifacts.md}.
\item Inspect the named frozen bundle files directly, starting with \texttt{manifest.json}, \texttt{aggregate\_table.md}, \texttt{run\_rows.jsonl}, and the significance/bootstrap files for the cited run.
\item Use \texttt{DATASET\_MANIFEST.md} to verify dataset provenance and frozen input status.
\end{itemize}

Together these materials make the empirical law auditable from the paper package itself. The manuscript therefore depends on frozen local evidence, not on untracked reruns or external archive links.
